\documentclass[letterpaper]{report}
\usepackage{makeidx,verbatim,texhelp,fancyhea}
\input{psbox.tex}
\parskip=10pt
\parindent=0pt
\title{JAZZ++ MIDI sequencer version 5.0}%
\author{Andreas Voss $<$Andreas.Voss@jazzware.com$>$\\
Per Sigmond $<$Per.Sigmond@jazzware.com$>$}%
\date{February 2008}%
\makeindex%
\begin{document}%
\maketitle%
\pagestyle{fancyplain}%
\tableofcontents%

%%%%%%%%%%%%%%%%%%%%%%%%%%%%%%%%%%%%%%%%%%%%%%%%%%%%%%%%%%%%%%%%%%%%%%%%%

\chapter{Introduction}\label{introduction}

JAZZ++ is a full size MIDI sequencer with audio support. In addition to basic
sequencer functions like record and play, JAZZ++ provides many edit features
like quantize, copy, transpose, graphical pitch editing, multiple
undo/redo etc.

JAZZ++ has two main windows; one
operating on the song as a whole, and one operating on single tracks and
single events. Among the other windows and dialogs are mixer windows for easy
adjusting of default track parameters, and all settings are saved into the
MIDI file for later use.

JAZZ++ has support for {\em GM} (General MIDI) equipment with both {\em GS}
and {\em XG} extensions.

JAZZ++ offers special functions like GS/XG sound editing, setting of effect
types, a {\em random rhythm generator} and a {\em harmony browser}. Windows
versions of JAZZ++ can be synchronized to other sound equipment (e.g. a tape
recorder) using MIDI Time Code (MTC) or SongPointer synchronization
techniques.

As an extension to normal MIDI sequencing, JAZZ++ provides audio play
(*.wav files), audio recording and audio/midi integration. A multi function
{\em audio editor} is also included.

\xlpignore{\begin{figure}%
\latexonly{\image{13cm;0cm}{jzwins.ps}}%
\htmlonly{\image{0cm;0cm}{jzwins.png}}%
\winhelponly{\image{0cm;0cm}{jzwins.bmp}}%
\caption{Main windows of JAZZ++}%
\end{figure}}

\section{System requirements}\label{requirements}

\subsection{MIDI and audio equipment}
To play MIDI music your PC must be attached to a MIDI synthesizer device. The
synthesizer can either be an integrated part of a sound-card, or an external
device attached to the PC with a cable.  To record MIDI music you will need a
midi-capable piano-keyboard attached to the PC (as an alternative you can
enter notes using the mouse).

For audio you need an audio capable sound card.

\subsection{Linux operating systems}

The device-independent part of JAZZ++ (the sequencer / editor part) is known
to run on Linux, Sun-OS4 and Solaris2 systems, but should compile and run
(with some minor hacking) on any system supporting Linux/X11. To compile you
will need the wxWidgets GUI package together with GTK.

The device-dependant part (midi-driver) currently supports:
\begin{itemize}
\item OSS/Free (OSS/Linux) midi driver API (formerly known as USS/Lite and
VoxWare). This API is included in later Linux kernels. The JAZZ++ code
interfacing OSS has only been tested on Linux platforms.
\item JAZZ++'s native MPU-401 midi driver (requires an intelligent-mode MPU-401
interface-card). The driver is implemented as a Linux kernel
load-module. In this mode JAZZ++ works as a TCP/IP-based client/server
application allowing the sequencer/editor and the midi-driver to run on
separate computers. The native driver has better support for the MPU-401
than the JAZZ++/OSS solution (e.g. support for external timing source).
\end{itemize}

\subsection{Windows operating systems}

JAZZ++ supports the windows MIDI / audio driver interface. A driver for your
MIDI interface card or sound card must be installed before JAZZ++ can play /
record MIDI or audio.

Provided such a MIDI driver is installed in your system, JAZZ++ currently
supports 32-bit platforms.

\subsection{CPU and memory}
For MIDI-only operation a 486 CPU (or equivalent) is sufficient in most
cases.  For satisfactory audio operation a Pentium CPU with 32 MBytes
RAM is recommended.


\section{Where to get it}\label{wheretoget}

For latest updates on software from {\em sourceforge}, please look at the
JAZZ++ homepage at

\begin{indented}{2cm}
{\tt http://jazzplusplus.sourceforge.net}
\end{indented}

You are welcome to join the JAZZ++ mailing list by sending mail to

\begin{indented}{2cm}
{\tt jazz-users-request@hia.no}
\end{indented}

with subject-field containing {\tt subscribe}.

We'd be happy to receive your bug reports (well, somehow) or success
stories. And - of course - you are welcome to contribute to the JAZZ++
project by sending code, suggestions, ..., or corrections for this
manual.



\section{Copyright}\label{copyright}

The JAZZ++ MIDI Sequencer is copyright (C) 1994-2000 Andreas Voss and
Per Sigmond, all rights reserved.

JAZZ++ program is free software; you can redistribute it and/or modify
it under the terms of the GNU General Public License as published by
the Free Software Foundation; either version 2 of the License, or
(at your option) any later version.

This program is distributed in the hope that it will be useful,
but WITHOUT ANY WARRANTY; without even the implied warranty of
MERCHANTABILITY or FITNESS FOR A PARTICULAR PURPOSE.  See the
GNU General Public License for more details.

You should have received a copy of the GNU General Public License
along with this program; if not, write to the Free Software
Foundation, Inc., 675 Mass Ave, Cambridge, MA 02139, USA.




\section{Installation and Startup}\label{installation}

For Windows there is an install utility (run {\tt setup.exe}) which
mostly automates the installation process.

For Linux there is a shell script {\tt install-jazz} in the
distribution. You must be logged in as root when running the installation
script. As a default JAZZ++ will be installed to {\tt /usr/local/jazz},
this may be changed by editing {\tt install-jazz}. Also you
must set the environment variable JAZZ to point the jazz directory by
adding the line:

{\tt export JAZZ=/usr/local/home}

to your .profile or .bashrc. Please refer to the file
{\tt install.txt} which contains up-to-date instructions.

On startup JAZZ++ looks for its configuration files
\begin{enumerate}
\item in the current directory
\item in the home directory (where the environment variable HOME points to)
\item in the directory where the environment variable JAZZ points to
\item in the directory from where JAZZ++ itself was started.
\end{enumerate}
The main configuration file is named {\tt jazz.cfg}.
This file also contains pointers to other configuration files {\tt (*.jzi)}
present in the same directory.
Then JAZZ++ searches for a song named {\tt jazz.mid} in the current directory
and loads it - if found. This song may contain default settings for your
midi equipment.


\subsection{Command line options (Unix only)}

Command line options (in any order):

\begin{itemize}
\item {\tt jazz -h}                                   -- Print help text
\item {\tt jazz -v}                                   -- Print help text
\item {\tt jazz -f filename}                          -- Input midi file
\item {\tt jazz -trackwin xpos ypos [width height]}   -- Position/size of main window
\item {\tt jazz -pianowin xpos ypos [width height]}   -- Position/size of piano roll window
\end{itemize}


\section{Acknowledgements}

JAZZ++ is based on the work of countless other people who have have made
their software freely available. Beside all the people who made Linux,
GNU and X11 possible, we'd like to mention

\begin{itemize}
\item {\em Hannu Savolainen} for writing all these excellent midi drivers, so
JAZZ++ can talk to nearly every midi card on Linux.
\item {\em Julian Smart} who wrote the free, portable GUI library
wxwin, so JAZZ++ runs with XView, Motif and MS-Windows.
\end{itemize}

We would also like to thank Henrik Seidel and Stefan Nitschke for code
contributions.


%%%%%%%%%%%%%%%%%%%%%%%%%%%%%%%%%%%%%%%%%%%%%%%%%%%%%%%%%%%%%%%%%%%%%%%%%
\chapter{Track window}\label{trackwin}

The trackwindow shows bars and tracks of your song while the
\helpref{pianowin}{pianowin} shows single events of one track. Events are shown like
bar codes, i.e. every event is shown as a short vertical line.

\xlpignore{\begin{figure}%
\latexonly{\image{13cm;0cm}{trackwin.ps}}%
\htmlonly{\image{0cm;0cm}{trackwin.png}}%
\winhelponly{\image{0cm;0cm}{trackwin.bmp}}%
\caption{Track window}%
\end{figure}}


\section{Toolbar}\label{twtoolbar}

From left to right there are the following buttons:
\begin{enumerate}
\item \helpref{Open a file}{file}
\item \helpref{Save actual song to file}{file}
\item Clear song (start a new one)
\item \helpref{Replicate}{replicate}
\item \helpref{Delete}{delete}
\item \helpref{Quantize}{quantize}
\item \helpref{Mixer}{mixer}
\item \helpref{Piano Window}{pianowin}
\item \helpref{Play}{recplay}
\item \helpref{Loop play}{recplay}
\item \helpref{Record}{recplay}
\item \helpref{Metronome}{metro}
\item Zoom in
\item Zoom out
\item \helpref{Undo}{undo}
\item \helpref{Redo}{undo}
\item \helpref{Panic (MIDI reset)}{panic}
\item \helpref{Help}{twhelp}
\end{enumerate}

\xlpignore{\begin{figure}%
\latexonly{\image{13cm;0cm}{twtoolb.ps}}%
\htmlonly{\image{0cm;0cm}{twtoolb.png}}%
\winhelponly{\image{0cm;0cm}{twtoolb.bmp}}%
\caption{Trackwin toolbar}%
\end{figure}}

\section{Selecting events}\label{select}

Many functions (e.g 'Delete' from the edit menu) require that you have
selected some range of bars and tracks that will be affected by the
selected operation. You can drag a rectangle to select a range. Holding
the shift-key down while pressing the left mouse button will continue a
selection. This way, you can select a range bigger than the window by
continuing a selection after scrolling. Alternatively, you can click
into the leftmost column to select a whole track.

The selection indicates a range of tracks (y-axis) and a range of bars (x-axis).
These values are copied to the Event filter from Settings-Menu. In the event filter
you can specify more precise, which events shall be worked with.

Example: To delete all events but program changes from bar 5 to 15 on track 1 to 3,
mark bars and tracks by dragging a rectangle. Then select the event filter and
disable program-changes (patch). Finally select Delete from Edit-Menu.

\xlpignore{\begin{figure}%
\latexonly{\image{13cm;0cm}{twselect.ps}}%
\htmlonly{\image{0cm;0cm}{twselect.png}}%
\winhelponly{\image{0cm;0cm}{twselect.bmp}}%
\caption{Selecting events}%
\end{figure}}

\section{Record and play}\label{recplay}

\xlpignore{\begin{figure}%
\latexonly{\image{13cm;0cm}{recplay.ps}}%
\htmlonly{\image{0cm;0cm}{recplay.png}}%
\winhelponly{\image{0cm;0cm}{recplay.bmp}}%
\caption{Record and play}%
\end{figure}}

There are 2 play modes
\begin{itemize}
\item {\em play} - normal linear play
\item {\em loop play} - requires that a loop-play range is selected (some bars in a track)
\end{itemize}

Each of the play modes has its own toolbar button. The buttons act both
as start-play and stop-play. After using JAZZ++ for a while you will
probably find it more convenient to start/stop play by clicking in the
top line of the Trackwindow. Play will then start at the bar where you click.
Loop play can also be done from the top-line by pressing shift-click.

Recording is done by pressing the record button after selecting
some bars in a track. The recorded events will be merged into the selection.
You can also start recording by clicking in the top-line of the Trackwindow
(some bars must be selected). If you start record with the right mouse button,
the track will be muted before play starts and existing events in the
selection will be replaced by the recorded events.

For a step-by-step tutorial on how to make a midi song from scratch, look at
\helpref{How to build a new midi song}{exmidisong}.

\section{Metronome}\label{metro}

The metronome button can be toggled on and off.
See also \helpref{Metronome settings}{metroset}.

\section{Speed adjustment}\label{speed}

A single click with left/right mouse button into the speed field will
increase / decrease the speed by one beat / minute. Keep the mouse button
down to make bigger changes. Hold down the shift key to change
values in steps of 10 bpm.

\xlpignore{\begin{figure}%
\latexonly{\image{13cm;0cm}{speed.ps}}%
\htmlonly{\image{0cm;0cm}{speed.png}}%
\winhelponly{\image{0cm;0cm}{speed.bmp}}%
\caption{Speed adjustment}%
\end{figure}}


\section{Track Defaults (program, volume etc)}\label{trackdefs}

Clicking at the top of the column named Prg cycles through the available
defaults (Prg = program change, Bnk = program bank (used in combination
with Prg for GS and XG sounds), Vol = volume, Pan = pan, Rev = reverb,
Cho = chorus). These events are sent before play starts. To adjust a
value, click left or right mouse button to increase or decrease the
value (use the shift key for quick increase or decrease). A value of 0
means there is no event to be sent. Please note that the defaults can also be
edited (graphically) in the \helpref{Mixer dialog}{mixer} selected from the
\helpref{mixer toolbar button}{twtoolbar} or the {\em Parts} menu.

Clicking the unnamed column containing P's toggles the Trackstatus:
P=play, S=solo, M=mute.

Clicking into the top of the leftmost column (labeled 'M') toggles
display of midi channel and track number.

\xlpignore{\begin{figure}%
\latexonly{\image{13cm;0cm}{trackdef.ps}}%
\htmlonly{\image{0cm;0cm}{trackdef.png}}%
\winhelponly{\image{0cm;0cm}{trackdef.bmp}}%
\caption{Track defaults}%
\end{figure}}

\section{Trackname, midi channel etc}\label{trackdlg}

Click on the name-field will open a Dialog. You can enter a new name for
the track, select the instrument to be played and define default midi channel.
If 'force all events to midi channel' is selected, all existing events on
the track will be set to the default midi channel when pressing 'Ok'.
Otherwise the midi channel will not affect events
already on the track.

\xlpignore{\begin{figure}%
\latexonly{\image{13cm;0cm}{tracknam.ps}}%
\htmlonly{\image{0cm;0cm}{tracknam.png}}%
\winhelponly{\image{0cm;0cm}{tracknam.bmp}}%
\caption{Track name, midi channel etc.}%
\end{figure}}


\section{Moving a whole track}

A whole track can be moved by clicking the right mouse button into the
track-name field, drag down/up to the destination and then release the
mouse button (drag and drop).


\section{File menu}\label{file}

\begin{itemize}
\item {\em Load, Save, Save As} - Load/Save MIDI-Standard files.
JAZZ++ will load both Format 0 and Format 1 midi files. Jazz always
saves Format 1 files.

\item {\em New} - Start editing a new song (clears the presently loaded
song). You will be asked for a file name when saving the new song.

\item {\em Load Template} is like normal load except that JAZZ++ will
treat the song as if it was a new song (you will be asked for a file
name when saved).

\item {\em Load Pattern} merges a midi standard file into the song. After
selecting the file, click with the left button to the destination position.
To abort loading after selecting the file click with the right mouse button
somewhere.

\item {\em Save Pattern} saves the selected events into a midi standard
file (Format 1).

\item {\em Quit} - exit JAZZ++.

\end{itemize}


\section{Edit menu}\label{twedit}


\subsection{Cut, Copy and Paste}\label{twclip}

This enables you to cut or copy some part of a song into an internal buffer
and paste it to another position of the song again. This can be used to
copy events from one song to another.

Do do this, \helpref{select}{select} some events and then select Cut or
Copy from the edit menu. To paste the events select paste from the edit menu
and click to the destination track/bar. To abort a paste operation click with
the right mouse button somewhere. Pasted events are merged with existing
events, so you may want to erase the destination tracks before pasting.



\subsection{Replicate}\label{replicate}

\xlpignore{\begin{figure}%
\latexonly{\image{13cm;0cm}{replic.ps}}%
\htmlonly{\image{0cm;0cm}{replic.png}}%
\winhelponly{\image{0cm;0cm}{replic.bmp}}%
\caption{Replicate dialog}%
\end{figure}}

To replicate (copy) some events, you must first \helpref{select}{select}
an area. Then, after you have invoked {\em Edit$->$Replicate} (or
pressed the \helpref{replicate button}{twtoolbar}), your cursor will change to
a '$=>$', meaning that you should mark a
destination for the selected area (clicking with the right mouse button will
abort replication). After clicking destination position, you get a dialog
box. The meanings of the different check boxes are:

\begin{itemize}
\item {\em Erase Destin}

\begin{itemize}
\item {\em Checked} - destination will be replaced by source
\item {\em Unchecked} - events will be merged
\end{itemize}

\item {\em Erase Source}. Source will be moved to destination

\item {\em Repeat Copy}. After clicking Ok, you can click at the end of the
destination range. Repeated copies of source will fill the destination range.

\item {\em Insert Space}. All Events right from the destination point will
be moved to the right to make room for the copied events.

\end{itemize}

Source and Destin may overlap.



\subsection{Delete}\label{delete}

\xlpignore{\begin{figure}%
\latexonly{\image{6cm;0cm}{delete.ps}}%
\htmlonly{\image{0cm;0cm}{delete.png}}%
\winhelponly{\image{0cm;0cm}{delete.bmp}}%
\caption{Delete dialog}%
\end{figure}}

To delete some events, you must first \helpref{select}{select} an area.
Then, after you have invoked {\em Edit$->Delete$} (or pressed the
\helpref{delete button}{twtoolbar}), you get a
dialog box. The check box 'Leave space' has the following meaning:

\begin{itemize}
\item {\em checked} - the deleted events leave a hole
\item {\em unchecked} - the events from right are moved to the left to
fill up the hole
\end{itemize}


\subsection{Quantize}\label{quantize}

Quantize will put Note-On events in the selected area 'in time' (moved to the
nearest step-timing value). Only Note-On events are changed.

\xlpignore{\begin{figure}%
\latexonly{\image{8cm;0cm}{quantize.ps}}%
\htmlonly{\image{0cm;0cm}{quantize.png}}%
\winhelponly{\image{0cm;0cm}{quantize.bmp}}%
\caption{Quantize dialog}%
\end{figure}}

To quantize some events, you must first \helpref{select}{select} an area.
Then, after you have invoked {\em Edit$->$Quantize} (or pressed the
\helpref{Quantize button}{twtoolbar}), you get a
dialog box. The 'steps' listbox represent the 'granularity' of the
quantization, and should normally correspond to the smallest note-timing
value of the music to be quantized. The checkboxes control whether the
start of the note or the end of the note (or both) should be quantized.

There is also a slider named 'Groove' that can be used for quantizing notes
'off beat'. A value of e.g. +20 means that the notes are moved to the nearest
step-timing plus 20 percent of a step-value (notes will sound a bit 'late'). A
negative value will make notes sound a bit 'early'. For normal quantization
the value should be 0.


\subsection{Set MIDI Channel}\label{setchan}

The selected events will be set to the selected MIDI Channel.

\xlpignore{\begin{figure}%
\latexonly{\image{8cm;0cm}{setchan.ps}}%
\htmlonly{\image{0cm;0cm}{setchan.png}}%
\winhelponly{\image{0cm;0cm}{setchan.bmp}}%
\caption{Set MIDI channel dialog}%
\end{figure}}


\subsection{Transpose}\label{transpose}

\xlpignore{\begin{figure}%
\latexonly{\image{8cm;0cm}{transpos.ps}}%
\htmlonly{\image{0cm;0cm}{transpos.png}}%
\winhelponly{\image{0cm;0cm}{transpos.bmp}}%
\caption{Transpose dialog}%
\end{figure}}

If no scale is selected, Amount is the number of semitones to transpose, e.g. an
Amount of 1 will change C to C\#, C\# to D etc. If the scale C is selected, C
will become D, D will become E, E will become F etc. JAZZ++ tries to analyze
the key of the selection and displays it in the dialog.

The scale 'selected' is the set of notes in the selection. If eg the selection
contains the notes C, D, E an amount of 1 will change C's to D's, D's to E's and
E's to C's. Its useful in the piano window create inversions of chords.

Fit into scale with an amount of 0 will change all notes to the nearest one
contained in the selected scale. If eg the selection contains some C\# and you
select C scale, all C\#'s will become C's (D would be possible too, but transpose
down is preferred in this case). If amount is not 0, the selection is transposed
in semitones first and 'Fit into scale' afterwards.

Only Note-On events are changed.


\subsection{Velocity}\label{veloc}

\xlpignore{\begin{figure}%
\latexonly{\image{8cm;0cm}{veloc.ps}}%
\htmlonly{\image{0cm;0cm}{veloc.png}}%
\winhelponly{\image{0cm;0cm}{veloc.bmp}}%
\caption{Set velocity dialog}%
\end{figure}}

Changes the velocity of Note-On events. If Stop is 0, all Note-On events will get
the Start velocity. If Stop is greater 0, events at the beginning of the selection
will get the Start velocity, those at the end will get the Stop velocity
(e.g. crescendo). The choice Add/Sub/Set determines, how the value is applied.


\subsection{Length}\label{length}

Adjust length of Note-On events. see \helpref{Velocity}{veloc}.


\subsection{Shift}\label{shift}

\xlpignore{\begin{figure}%
\latexonly{\image{8cm;0cm}{shift.ps}}%
\htmlonly{\image{0cm;0cm}{shift.png}}%
\winhelponly{\image{0cm;0cm}{shift.bmp}}%
\caption{Shift dialog}%
\end{figure}}

Moves events left or right in amounts smaller than a bar. The snap value is
adjusted in the pianowin.

\subsection{Cleanup}\label{cleanup}

\xlpignore{\begin{figure}%
\latexonly{\image{6cm;0cm}{cleanup.ps}}%
\htmlonly{\image{0cm;0cm}{cleanup.png}}%
\winhelponly{\image{0cm;0cm}{cleanup.bmp}}%
\caption{Cleanup dialog}%
\end{figure}}

Deletes note-on events that are shorter than a selected length. Optionally
shorten notes to insure that notes of the same pitch does not overlap.

Accidental keyboard hits often leave very short notes on the track that
distorts the music. Another problem is that two overlapping notes of the
same pitch causes the second note to stop sounding (because of the note-off
event from the first note). This cleanup function effectively solve such
problems.

\subsection{Search Replace}\label{search}

\xlpignore{\begin{figure}%
\latexonly{\image{13cm;0cm}{search.ps}}%
\htmlonly{\image{0cm;0cm}{search.png}}%
\winhelponly{\image{0cm;0cm}{search.bmp}}%
\caption{Search Replace dialog}%
\end{figure}}

Searches for controller events and changes the controller number, eg
transforms modulation wheel to volume.




%%%%%%%%%%%%%%%%%%%%%%%%%%%%%%%%%%%%%%%%%%%%%%%%%%%%%%%%%%%%%%%%%%%%%%%%%

\section{Parts menu}\label{partsmenu}

This menu deals with all kinds of instrument settings, or {\em parts} settings.
Some of the parameters will work with
almost any General Midi (GM) equipment (like Volume settings) while
others only work with GS- or XG-compatible synthesizers. The number of
selectable menu entries will vary depending on the
\helpref{synthesizer type}{synthtype} settings.

Please note that some
menu entries may have no effect on your equipment even if selectable. In some
cases special {\em system exclusive} message needs to be sent
first in order to enable the specific feature
(see \helpref{Edit/create single events}{pweventedit} on how to create a
system exclusive event).
In other cases the equipment just doesn't support the feature.
Please consult the owner's manual of your sound generator for more info.

All parameters are operated with sliders per track (part). Note that for some
platforms the number of adjustable parameters are restricted because of
the size of the dialogs.

Please note that only one {\em parts} dialog can be opened at a time.

\subsection{Mixer and Master}\label{mixer}

Setting of default {\em Volume}, {\em Pan}, {\em Reverb} and {\em Chorus}.

\xlpignore{\begin{figure}%
\latexonly{\image{13cm;0cm}{mixer.ps}}%
\htmlonly{\image{0cm;0cm}{mixer.png}}%
\winhelponly{\image{0cm;0cm}{mixer.bmp}}%
\caption{Mixer window}%
\end{figure}}

\subsection{Sound, Vibrato and Envelope}

Sound dialog - change sound color:

\begin{itemize}
\item {\em Cutoff frequency} - filter high frequencies
\item {\em Resonance} - make the sound more 'synthetic' (wha-wha sound)
\end{itemize}

Vibrato dialog - change the way the tone's vibrato develops in time:

\begin{itemize}
\item {\em Vibrato Rate} - controls how fast the vibrato will be
\item {\em Vibrato Depth} - controls how much vibrato you get
\item {\em Vibrato Delay} - controls at what time after note start the
vibrato begins
\end{itemize}

Envelope dialog - change the way the tone envelope (amplitude) develops in
time:

\begin{itemize}
\item {\em Envelope Attack} - controls how fast the tone increases in volume
when it starts
\item {\em Envelope Decay} - controls how fast the tone will
decrease in volume after it has started (also called {\em sustain})
\item {\em Envelope Release} - controls how fast the tone decreases in volume
after it is has ended (how fast the tone stops)
\end{itemize}

\subsection{Controller dialogs}
These dialogs change the way the different controllers affect the sound.

Controllers are:

\begin{itemize}
\item {\em Pitch Bender} - used to control continuous tone-pitch in real time
\item {\em Modulation Wheel} - used to control 'tone expression'
(e.g. vibrato) in real time
\item {\em Channel Aftertouch (CAf)} - some keyboards senses how hard you
press down the key after the attack, used for 'tone expression' in real
time. Affects all tones of the part (midi channel).
\item {\em Polyphonic Aftertouch (PAf)} - same as above, but only affects
one key pitch of a part (channel).
\item {\em Continuous Controller 1 (CC1)} - used for misc. 'tone expression'
\item {\em Continuous Controller 2 (CC1)} - used for misc. 'tone expression'
\end{itemize}

All the controllers have three associated types dialogs with sliders;
{\em Basic}, {\em LFO1} and {\em LFO2}.
The settings tell how much the controller will trigger the
sound-change described by each parameter. Parameters are:

\begin{itemize}
\item {\em Pitch Control / Sense} - pitch in half-tone steps
\item {\em TVF Cutoff} - controls high frequency cutoff (TVF =
Time Variant Filter)
\item {\em Amplitude} - controls amplitude (dynamic volume)
\item {\em LFO1 Rate, Pitch, TVF and TVA} - controls the rate of the Low
Frequency Oscillator 1 (LFO1) and how it affects pitch, TVF and TVA (TVA =
Time Variant Amplitude)
\item {\em LFO2 Rate, Pitch, TVF and TVA} - controls the rate of the Low
Frequency Oscillator 2 (LFO2) and how it affects pitch, TVF and TVA (TVA =
Time Variant Amplitude)
\end{itemize}

The {\em Basic} dialogs of the {\em Continuous Controller} entries also has
a {\em Controller Nr} dialog to set what continuous controller the part
will respond to.


\subsection{Drum Instrument Parameters}

This is a dialog for controlling different parameters of the individual drum
instruments of the primary drum part. The track selected is the first track
with channel equal to the {\em .drumchannel} setting (in the configuration
file ). Note that the name of the track also has to be set.

Drum instruments are listed in a list box. The sliders are valid for the
instrument currently selected from the list. Listed instruments to be
controlled are added/deleted by using the {\em Add} and {\em Del} buttons.

The controllable parameters are:

\begin{itemize}
\item Drum instrument {\em pitch coarse}
\item Drum instrument {\em TVA} (volume)
\item Drum instrument {\em panpot} (checkbox gives random pan)
\item Drum instrument {\em reverb send}
\item Drum instrument {\em chorus send}
\end{itemize}


\subsection{Partial Reserve}

Sound generators have a limited amount of tones that can be played
simultaneously. This is sometimes referred to as 'maximum polyphony'.
When the maximum polyphony is exceeded the synthesizer will start
cutting tones, which can cause audible 'distortion' to the music.

In some GS equipment you can assign the number of tones (oscillators)
that is assigned to each part (channel) and thereby control which instruments
are 'more important'. This is what this dialog is for.

\subsection{Part Mode}

Here you can set two things:

\begin{itemize}
\item {\em Rx Channel} - the channel to be received by the part
\item {\em Use for rhythm} - sets the part to normal mode (value 0)
or drum modes (value > 1). Part 10 (channel 10) has default value 1,
all others have default value 0. On GS and XG synths you can get another drum
part by setting that part to value > 1.
\end{itemize}

%%%%%%%%%%%%%%%%%%%%%%%%%%%%%%%%%%%%%%%%%%%%%%%%%%%%%%%%%%%%%%%%%%%%%%%%%
\section{Settings menu}\label{twsettings}

\subsection{Filter}\label{midifilter}

Using the 'event filter dialog', you can specify what types and ranges of
events that will be affected by \helpref{selection}{select} operations like
e.g. 'Delete' and 'Replicate'.

\xlpignore{\begin{figure}%
\latexonly{\image{13cm;0cm}{filter.ps}}%
\htmlonly{\image{0cm;0cm}{filter.png}}%
\winhelponly{\image{0cm;0cm}{filter.bmp}}%
\caption{Event filter dialog}%
\end{figure}}

The entries {\em From Time}, {\em To Time}, {\em From Track} and
{\em To Track} are automatically set by selecting a range with the mouse.

The checkboxes {\em Note}, {\em Controller}, {\em Patch} (Program Change),
{\em Pitch}, {\em Meter} and {\em Other} decide whether these event types
should be included in the selection. For the first four event types you can
also specify more closely which values will be included.

By default, all 'boxes' except {\em Other} are checked, leaving e.g.
{\em Tempo} events unchanged by selection operations.

\subsection{Window}
adjust size of fonts (y-axis) and events (x-axis).



\subsection{Song}\label{songset}

\xlpignore{\begin{figure}%
\latexonly{\image{10cm;0cm}{songset.ps}}%
\htmlonly{\image{0cm;0cm}{songset.png}}%
\winhelponly{\image{0cm;0cm}{songset.bmp}}%
\caption{Song settings dialog}%
\end{figure}}

Here you can adjust the overall resolution of the song in ticks per quarter.
The native JAZZ++ driver (MPU-401) supports distinct values only. With OSS
and on MS-Windows you can use any value here. If you change this value,
all event timings are recomputed.

A big value of 'ticks per quarter' means that the time resolution of the
notes in the song will be better, but the strain on the sequencer system
will also increase. Many people think it is important to have a big value here,
while others say you won't hear the difference. We recommend you to use
the default value (120 ticks per quarter) which gives sufficient resolution
for any 'normal' midi song.

In this dialog you may also set the maximum length of the song in bars.
This value is used to set the scrollbars in the track and piano windows.


\subsection{Metronome Settings}\label{metroset}

\xlpignore{\begin{figure}%
\latexonly{\image{10cm;0cm}{metroset.ps}}%
\htmlonly{\image{0cm;0cm}{metroset.png}}%
\winhelponly{\image{0cm;0cm}{metroset.bmp}}%
\caption{Metronome settings dialog}%
\end{figure}}

Adjust metronome instruments and velocity.

An ``accented'' metronome means that the first count of a bar sounds different
than the other counts.

\subsection{Effects}

Adjust GS or XG effect settings. In XG mode, also equalizer settings can be
chosen.

In GS mode, effects can be adjusted in two modes dependant on the settings
of the 'Use effect macro' and 'Use chorus macro' checkboxes:
\begin{itemize}
\item {\em macro mode} - the choices in the macro listboxes are valid
\item {\em individual parameter mode} - the values of the parameter
sliders are valid
\end{itemize}

Macro mode will be sufficient in most practical cases. Setting of individual
parameters is more difficult, but can give very good results when done right.

The mode settings (position of the checkboxes) can be saved with the
"Save settings" menu entry. The current values of the listboxes or parameter
sliders are stored in the midi file.

\subsection{Timing}

Controls clock source and real-time output.

The available choices of {\em clock source} will vary depending on system
platform:

\begin{itemize}

\item {\em Internal}. The sequencer is itself the source of the clock timing
(available on all system platforms).

\item {\em Song Pointer}. The sequencer clock locks to midi real-time messages
received from midi input. The source of these messages is typically a
tape-machine with an external sync-box or just another sequencer. Song
pointer sync is available on Unix with JAZZ++-native MPU-401 driver and on
MS-Windows.

\item {\em Midi Time Code}. The sequencer clock locks to Midi Time Code
(MTC) messages received from midi input. The source is typically a
tape-machine with an external sync-box doing SMPTE-to-MTC conversion or
a MTC-capable recording machine. MTC sync is available on MS-Windows platform.

\item {\em FSK}. The sequencer clock locks to FSK information received from
tape. This is available on Unix with an MPU-401 card supporting FSK
sync.

\end{itemize}

How to sync?

\begin{itemize}

\item {\em Internal}.
Just play; JAZZ++ is the master and decides all timing
using internal clocks in the computer.

\item {\em Song Pointer}.
Make sure the device (here called "tape") you are syncing
to sends midi real-time messages (songpointer, midi clock etc.) This could
involve recording a "songpointer stripe" to a tape track by means of
a sync box. To do this you must play the song with the sequencer in
"Internal" sync mode and with "Realtime to MIDI Out" enabled. (For more info
refer to the instructions of the sync box.)

To play, enable "Song pointer" sync mode in sequencer (JAZZ++), start play on
sequencer
first, then start tape at any point in the song. The sequencer is now locked
to the movements of the tape. (Note that some sync boxes don't send
song-pointer at beginning of the song. In this situation you must manually
start sequencer at song-start before running tape.)

\item {\em Midi Time Code}.
Make sure the device (tape) you are syncing to sends
MTC messages (could involve recording an SMPTE stripe). Set the MTC offset
and the MTC type either manually or by recording it while listening to the
tape. (The MTC offset is the point in time when the sequencer is supposed
to start playing the song.) When recording MTC offset the sequencer will
automatically be set into MTC sync mode.

To play, set sequencer to MTC sync mode, start play on sequencer, then start
tape at any point in the song. The sequencer is now locked to the movements
of the external device.

\item {\em FSK}.
An "FSK stripe" must first be recorded to a tape track. To do
this you must play the song with the sequencer in "Internal" sync mode. Be
sure to start FSK stripe recording some time before sequencer is started
(must record the pilot tone).

To play, set sequencer to FSK sync mode, start tape in the pilot-tone area,
then start sequencer before pilot tone ends. (Note that with FSK play must
always start at beginning of the song.)

\end{itemize}

The "Realtime to Midi Out" choice is useful for syncing another sequencer
to this one or to record a "songpointer stripe" to tape via a sync box.
Note that "Realtime to Midi Out" will only work properly for a
ticks-per-quarter resolution dividable by 24
(see \helpref{Song settings}{songset}).


\subsection{Midi Thru}
Enable/disable midi thru. Should normally be enabled.

When enabled, the sequencer will send all notes received on
the MIDI IN port to the MIDI OUT port. This basically means that you will hear
what you play on the keyboard.

\subsection{Synthesizer Type Settings}\label{synthtype}

\xlpignore{\begin{figure}%
\latexonly{\image{13cm;0cm}{syntype.ps}}%
\htmlonly{\image{0cm;0cm}{syntype.png}}%
\winhelponly{\image{0cm;0cm}{syntype.bmp}}%
\caption{Setting of synthesizer type}%
\end{figure}}

Many MIDI sound generators (a synthesizer, soundmodule or soundcard) offer
possibilities to choose an extended number of instruments (more than
the 128 General MIDI instruments), or to control things like sound colour,
reverb/chorus settings or time-variant filters. This is done by sending
specific MIDI commands from the sequencer
(bank-, NRPN- and SYSEX-events). However, the parameters of
these commands are different depending on which standard your sound
generator supports.

To allow JAZZ++ to get the most out of your MIDI sound generator, you should
choose the synthesizer type option that corresponds most closely to your
equipment.

Options are:

\begin{itemize}

\item{\em GM}
- Choose this if your sound generator is a standard General MIDI device.
This option is safe and will work for almost any sound generator.

\item{\em GS}
- GS is an extension to the GM specification that allows for 317 different
instruments, effect type control and a lot more. Look for the GS logo on
your equipment. GS was originally defined by Roland (R).

\item{\em XG}
- XG is another extension to GM. A lot of extra instruments
are defined as well as a number of extended control possibilities. Look
for the XG logo on your equipment. XG was originally defined by
Yamaha (R).

\item{\em Other}
- This option is for sound generators that do not follow any of the standards
mentioned above. JAZZ++ allows user customization of parameters like
instrument lists (patches), controllers and drum instrument names. Look
into the JAZZ++ \helpref{configuration files}{installation} for more info.

\end{itemize}

The 'Auto send MIDI reset' options control the conditions for when JAZZ++
should send a MIDI reset command (corresponding to the synthesizer type) to the
synthesizer device.

\begin{itemize}
\item {\em Never} - MIDI reset is never automatically sent to the device
(only when you press the \helpref{panic putton}{twtoolbar})
\item {\em Song start} - MIDI reset is sent whenever play starts from the
beginning of the song (default)
\item {\em Start play} - MIDI reset is always sent when play starts
\end{itemize}

Files saved from JAZZ++ will always contain a MIDI reset near the beginning
of track 0.

\subsection{Midi Device}

You may select the input and output devices from the list of installed
drivers.

\subsection{Save Settings}
Save settings to config file.
Affected settings: midi thru, clock source, realtime out, effect macro modes,
display welcome message and window geometry.


\section{Misc menu}

\subsection{Undo and redo}\label{undo}

Undo or redo the last operation. The last 20 operations can be undone/redone.


\subsection{Merge / Split Tracks}

{\em Split tracks} places all events from track 0 to track 1..16 depending on
the midi channel (useful for 'decoding' a format 0 midi file).

{\em Merge tracks} moves all events to track 0 (not so useful as JAZZ++ only
saves file format 1).

\subsection{Meterchange}\label{meterchange}

\xlpignore{\begin{figure}%
\latexonly{\image{6cm;0cm}{meter.ps}}%
\htmlonly{\image{0cm;0cm}{meter.png}}%
\winhelponly{\image{0cm;0cm}{meter.bmp}}%
\caption{Meter change dialog}%
\end{figure}}

Changes the meter (counts per bar and note-value of each count) starting
from specified bar to the end of song or to the
next meterchange. {\em Numinator} is the 'counts per bar', while
{\em Denominator} is the 'note-value of a count'.

Meterchange events are placed on track 0 always.

\subsection{Midi reset}\label{panic}

Sends an initialize message (either 'GM ON', 'GS ON' or 'XG ON') to the midi
port. If \helpref{synthesizer type}{synthtype} is correctly set, most sound
generators will initialize to default values when such a message is
sent.

\subsection{Harmony browser}\label{hbrowser}

\latexonly{\begin{figure}%
\image{13cm;0cm}{hb.ps}%
\caption{Harmony browser}%
\end{figure}}

\xlpignore{\begin{figure}%
\latexonly{\image{13cm;0cm}{hbtoolb.ps}}%
\htmlonly{\image{0cm;0cm}{hbtoolb.png}}%
\winhelponly{\image{0cm;0cm}{hbtoolb.bmp}}%
\caption{Harmony browser toolbar}%
\end{figure}}

The harmony browser shows the 4-tuned JAZZ++ harmonies of the four scales
types major, harmonic minor, melodic minor and harmonic major.

The first thing you can do is to click on chords and listen how they sound (eg
examples from a song book or a book about harmony theory). The selected
chord is placed into the pianowin copy buffer and can be pasted into
your song (see \helpref{mouse actions}{mouseact}.

The harmony browser is also able to analyze a section of your song and
to transpose parts of your song into new scales and chords. Also you may generate
chords thru the random rhythm generator into your song.

What you can do with the mouse:
\begin{itemize}
\item {\em left click on a chord} selects the chord and marks related chords (see below).
\item {\em shift + left} click selects the chord and appends it to the chord sequence
you need to transpose parts of your song
\item {\em right click} just listen to the chord, don't select it
\item {\em right click on the last chord} of the chord sequence removes this
chord from the sequence.
\end{itemize}

There is a toolbar with the following icons (from left to right):

\begin{itemize}

\item {\em The left four buttons} switch the scale type for which the standard
chords are displayed: major, harmonic minor, melodic minor and harmonic
major which is also known as ionic b13.

\end{itemize}

The following five Icons (with the cyan right arrow and the piano)
specify what related chords are marked when you select a chord:

\begin{itemize}

\item {\em 4,3,2,1,0} marks all chords that have 4,3,2,1,0 common notes with the
actual chord

\item {\em 1/2} marks all chords, that differ from the selected chord by exactly
one half note

\item {\em 251} marks all chords, that would come next in the 2-5-1 move

\item {\em =B} marks all chords, that share the same bass note with the selected chord

\item {\em Tri} marks all chords, that could be tritone substitutes (their bass note
is at a tritone interval from the selected chord).

\item {\em (piano)} mark all chords, that contain all of the keys in the pianowin
copy buffer. So you could select a few notes (max 4) in the pianowin and select
copy from the edit menu. Now the harmony browser will mark all chords that contain
those notes.

If the piano button is not selected, the harmony browser works the other way round.
Whenever you select a chord, the chord is placed into the pianowin buffer so you
can paste it into your song. The parameters (midi channel, length etc) can be set
in the midi-entry of the harmony browser menu.

\end{itemize}

The last four buttons do different things:

{\bf Haunschild layout} shows harmonies and scales in an order proposed
in a book about harmony theory by the German author Haunschild. If selected,
chords are shown in a way that 'valid moves' are

\begin{itemize}
\item {\em from left to right}, e.g the 2-5-1 move is displayed as three neighbored
chords. In general, all horizontal moves are allowed.
\item {\em vertical moves} are allowed, they change one semitone of a chord (at most).
\item {\em diagonal} from top-left to bottom right. I don't know why, but it sounds
nice.
\end{itemize}

{\bf Transpose} lets you change parts of your song to the harmonies and chords
of the chord sequence. When you select some bars in the trackwin and
click on transpose, then the selection is mapped to the chord sequence.
The default is, that on chord lasts one bar, but you may change this in the
settings menu in the harmony browser. If the results are not satisfying, you
should try to transpose track by track, this often works better. All in all,
transposing works better, when the selected material is not too complicated.

I had nice results with the following procedure: I recorded (very slow, as I am
not a good piano player) cj7 chords on one track, some monphone bass line on the
second track and a solo on the third track. Everything was in c major (white keys
only). Then I selected some chords and transposed each track on its own.

{\bf Analyze} Tries to find the chords and scales from the selection in the
trackwin. Again, the default is one harmony per bar which can be changed in the
settings menu. The algorithm looks at the length of the notes and chooses the
longest notes to be the notes of the chord.

{\bf Clear chord sequence} removes all chords from the chord sequence.


{\bf Menu entries}

\begin{itemize}
\item {\em edit chord} - often you may want to have chords/scales that are not displayed
in the harmony browser. You can select a chord from the chord sequence and change
it in this dialog.

\item {\em settings} - adjust parameters for transpose and analyze as explained above.
A value of 0 quarters per chord when transposing will map the chords to the length
of the selection, e.g if you have selected 16 Bars and your sequence has 4 chords,
every chord will last 4 bars. If you had chosen 4 quarters instead, the chord
sequence would have been repeated 4 times to fill the 16 bars of the trackwin selection.

\item {\em midi} - these settings affect what you hear, when clicking on a chord and what
is put into the pianowin buffer.

\end{itemize}

{\bf Generating chords with the random rhythm generator}

After defining a chord sequence, you may open the random rhythm
generator. When adding an instrument, it will show (beside drums) the
instrument 'chords from harmony browser'. You select it and define a more or
less randomized rhythm and proceed as described \helpref{here}{randrhy}.

\subsubsection{Transpose Algorithm}

This one is somewhat difficult to explain, but its probably the most
interesting feature of the Harmony Browser.  Lets assume, you have a
sequence of 4 chords. Then you select 4 bars in the track window (all
tracks) and you press the transpose button.  Now, the Harmony Browser
tries to change the harmonies and scales from the trackwindow selection
into the selected chords.

The algorithm looks at the first bar from the trackwin selection and
sums up the length of all different notes. It assumes, that the 4 notes with the
largest sum of length are the chords present in the song and changes them
to the 4 notes from the first chord of your chord selection. Afterwards,
the remaining notes from the trackwin selection are mapped to the remaining
notes of the scale the chord was taken from.



\subsection{Random rhythm generator}\label{randrhy}

\latexonly{\begin{figure}%
\image{13cm;0cm}{rrg.ps}%
\caption{Random rhythm generator}%
\end{figure}}

\xlpignore{\begin{figure}%
\latexonly{\image{13cm;0cm}{rrgtoolb.ps}}%
\htmlonly{\image{0cm;0cm}{rrgtoolb.png}}%
\winhelponly{\image{0cm;0cm}{rrgtoolb.bmp}}%
\caption{Random rhythm generator toolbar}%
\end{figure}}

This dialog generates rhythms from some kind of statistical specification.
You define probabilities for timing, velocity and note-length and from these
the rhythm is generated.

With the three sliders you specify the meter and (with \# Bars) how many
probabilities you want to define. Next, you can add or remove instruments
to the listbox. Then you specify the probabilities for rhythm, velocity
and length for each instrument.

Probabilities are set with the mouse where

\begin{itemize}
\item {\em left button} - draw some kind of curve
\item {\em right button} - clear values to zero
\item {\em + shift} - adjust all values simultaneously
\item {\em + ctrl} - adjust exactly one value
\end{itemize}

In the rhythm field you specify the probabilities for every time
position. A high value means a high probability that the instrument is
played definitely at that time. A low value means, that the instrument
{\em may} be played at that time. By adjusting high values only, you
can specify a definite rhythm without any randomization.

The velocity (and length) values are looked up by the generator, whenever
it decides to play a beat.  It chooses one value from 1..127 (x-axis,
1..8 for length) where positions with high probabilities are chosen more
often.

If the chosen length is greater than 1 step, there are no more events
generated until the note quits, even if there are high probabilities in the
rhythm field. You can abuse this, to make sure that there are not too many
events generated (example: you have rhythm probabilities at step 4 and 5,
but you want exactly {\em one} beat to be played at either of these.
Specifying a length of 2 will do this,  if step 4 is chosen, there cannot
be an event on step 5). The length may also be used for rhythms where you
want to have random intervals between the beats but you are not interested
in the absolute timing positions of the beats. Look at the cowbell definition
in the rrg1.rhy example for this.

We had the best results with the following: define every instrument
twice. In the first definition, select high probabilities on a few
timing positions together with high velocities, this makes the base
groove - it is loud and does not change very much. In the second
definition, have little probabilities on many timing positions with
a low velocity, this makes randomized background fills, which make the
whole thing more interesting and that do not override the base groove
because of the low velocity.

Aside from drums there are some special 'Instruments' you can select:

\begin{itemize}
\item {\em Controller} - opens another dialog to select a controller.
Randomizing works the same like with drum keys except that instead of
velocity the controller value is generated. With this, you may generate for
example random panpot events that move the sound between left and
right speaker or, in conjunction with the \helpref{Parts Menu}{partsmenu},
you can generate many different randomized sound effects.
\item {\em Pianowin all} - places the events currently selected in the piano
window. You could eg select a chord an then replicate it on random timing
positions.
\item {\em Pianowin one} - like pianowin all, but on every timing position
there is chosen one of all the selected events.
%\item[Harmony: chord, harmony: bass] only available, if the
%\helpref{harmony browser}{hbrowser} is opened and a chord sequence is defined. Places
%the selected chords or bass note randomly.
\end{itemize}

The events selected in pianowin are copied to an internal rhythm buffer, so
you may define multiple rhythms with different pianowin selections.

To generate, you select some bars on one track in the trackwin and press the
gen button.



\subsubsection{Instruments without randomizing}

If the checkbox 'randomize' is not selected, the velocity and length
fields are disabled. The rhythm is generated from the rhythm field only.
Every position with a bar height greater 0 gives a beat, the height of
the bar is the velocity.


\subsubsection{Instrument Groups}

Up to now every instrument has its own probability definition and plays
on its own without listening to the other instruments. By defining groups,
instruments can play a little more 'together'. Some instruments contribute
their rhythm to a group and others listen to the group and play
accordingly.

The contributors contribute their actual rhythm to the group, so the group
maintains its own rhythm that is something like the sum of the rhythms of
all contributing instruments.

The listening instruments modify their probabilities according to the
rhythm of the group. If the listen value is greater 0, then the instrument
will play at the same timing positions as the group rhythm. If the value
is negative, the instrument will play at the timing positions where the
contributors do not play.

The algorithm evaluates the instruments in the order they appear in the
list. So contributors should be defined before listeners so its actual rhythm
will be known to the listener. If a contributor is
defined after a listener, the listener will listen to what the contributor
played in the previous bar. An instrument may be listener and contributor
at once.

Example 1: There are three different conga instruments: muted,  high and low
and you want them to play exclusively (only one instrument at a time). The
first in the instrument list (eg the muted conga) will play its rhythm and
contribute it to a group. The next instrument (e.g. high conga) will listen
to the group with a value of -100, so it will play on those positions only
where the muted conga does not play. It will also contribute its
rhythm to the group. The third instrument (e.g. low conga) will also listen
with a value of -100 so it will play only on positions, where none of the
others play. See rrg2.rhy for an example of this.

Example 2: The open hi-hat shall go on some of the positions where the
bass drum plays, it shall not play alone (without bass drum). In this
situation the bass drum will contribute its rhythm to a group and the
open hi-hat will listen to the group with a value of +100.



\subsection{Mapper}\label{mapper}

This dialog allows to modify velocity, note length, note pitch and note start.
We'll explain this dialog by an example: Select length as 'Source' and velocity
as 'Destin' and select some tracks in the trackwin. In the lower part of the
dialog you can paint some kind of curve now, see section \helpref{Random
Rhythm}{randrhy} for details.

When you press the Apply-Button
the following happens: for every note on event, JAZZ++ looks at its length and
looks up the length on the x axis of the curve. Now it looks at the curve and
finds the corresponding value at the y-axis. This value is assigned to the
velocity of the note on event. So, if you draw a diagonal line from bottom
left to top right, short notes will get a low velocity and long note will get
a high velocity.

For rhythm this is analogous: the mapper looks up the start time in the x axis
and takes the value from the height of the corresponding bar.

When using random as 'Source' things are a bit different. The y-axis shows
probability values now and the x-axis shows possible Destin values.
JAZZ++ chooses for every note-on a new destin value - one of all the values
shown on the x-axis. But those values, with a high probability are chosen
more often, values with probability 0 will never be chosen.

The switch 'Add value' means, that the Destin value is added to the existing
value and the value is not replaced. You can subtract values too by adding
negative Destin values.

The rhythm and random features are probably the most interesting. The rhythm
allows, to adjust velocity to a certain groove and random allows to 'humanize'
midi events, add some very small timing inaccuracy or velocity changes.



\subsection{Random Shuffle}\label{shuffle}


With this dialog you can rearrange parts of your song based on random. The
selected bars are divided into small segments (default length is 1/4 note).
When pressing OK, jazz rearranges these segments in random order. Also,
some of the segments may be choosen more than once while other segments do
not appear at all. Options are:

\begin{itemize}

\item {\em Segment:} You can choose the length of the segments here. You may
choose very small segments (e.g. 1/16) to create new drum rhythms or large
segments (e.g 4/1) to rearrange large parts of the song.

\item {\em skip silence:} when checked, segments that contain silence (no notes)
are ignored. So you could place just a few events on some tracks in the piano
window and then replicate them in random order.

\item {\em random order:} when checked, jazz chooses the segments by random. If
not checked, jazz chooses the segments one after the other from left to right
and starts over on the left again when it reaches the right end. This makes
sense when skip silence is enabled and/or track mode is set to
'One track at a time'.

\item {\em track mode:} When you have selected more than one track, this
option applies. There are three choices:

\begin{itemize}

\item {\em All tracks in sync} means, that whenever jazz chooses a segment,
it will take the same segment from all tracks. For example if jazz chooses
the 7-th segment on the first selected track, it will place the 7-th segment
of all tracks.

\item {\em All tracks independently} means, that jazz chooses the segments
independently on every track. In the example, jazz may choose the 7-th
segment on the first track, the 3-rd segment on the second track and so on
and all these segments will sound together in the result.

\item {\em one track at a time} means, jazz will choose only one
segment from one track per step. If you have selected a drum track and
a percussion track, in the result there will be playing either drums or
percussion but not both together.

\end{itemize}

\end{itemize}


See the song/shuffle directory for an example. Load shuffle0.mid, then mark
from 5 to 24 on tracks 5 to 8 (chords, ..., muted gt). Select Shuffle from
the misc menu. Set the segment size to 1/4 note and set the trackmode to 'one
track at a time'. Then press ok. This should give you something similar to
shuffle1.mid. The song shuffle2.mid was generated from shuffle0.mid by deleting
the tracks 5 and 7 (marimba and brass). To delete a track, click on its
name and check the field 'Clear track' at the bottom of the track dialog.
After clicking OK, you can drag the empty track to a new position with
the right mouse button. To get shuffle2.mid set the segments size 1/8 note.


\subsection{Random Melody Generator}\label{shuffle}

With this dialog you can generate melodies, bass lines or chords.
In general it works as follows: first it generates a small piece of melody,
the seed. Then this seed is repeated several times and each time some variations are
applied to the seed. Finally the melody is mapped to the harmonies using the Harmony
Browser. See the description of the \helpref{Random Rhythm Generator}{randrhy}
for additional information. Adjustables:

\begin{itemize}

\item {\em steps/count, count/bar, \#bars} the maximum length of a seed

\item{\em pitch center, pitch range} defines the range of midi keys for the
generated melody.

\item{\em \#notes} one for melody and bass, two or three for chords

\item{\em random seed length} if checked, the length of the seed is
determined by random. If not, the seed length equals the length of
the rhythm field (see below) and produces more steady results for bass
and chords.

\item{\em rhythm as intervals} if checked, the rhythm field is interpreted as
intervals, if not, its interpreted absolute. Example: lets assume there is one bar
at 1 and one bar at 2. Interpreted absolute means, that notes may occur on the first
beat and/or on the second beat, but no notes will be generated at the third and fourth
beat. Interpreted as intervals means, that the time distance between two notes
may be either 1/16 or 4/16.

\item{\em note length} the length of a generated note. Like the Random Rhythm
Generator, a new note may start only after the old one has finished, so the
note length overrides the rhythm setting in concurrent situation.

\item{\em velocity} velocity setting, same as in Random Rhythm generation

\item{\em pitch interval} When the RMG generates the seed, it starts with a random
choosen pitch. Then for every next note, it adds an interval to its previous note.
The interval is choosen from the probabilties you define here. If, for example, there
is only one bar at 1, then every note will have an interval of 1 to its previous
note.

\item {\em seed length} Here you define the probabilities for the length of the seed.

\item {\em seed repeat} defines the probabilities for how often a sees will be
repeated

\item {\em seed variations} defines the variations that may be applied:
\begin{itemize}
\item {\bf 1} plain repeat with no variation
\item {\bf 2} multiplies the intervals by -1
\item {\bf 3} appends a new pattern to the beginning and to the end of the seed
\item {\bf 4} multiplies the intervals by 2 or 1/2
\end{itemize}

\item {\em rhythm} defines the probabilities for the rhythm.

\end{itemize}

Please take a look at the examples in the song/rmg directory. To generate a new
melody for example 1, do the following:
\begin{itemize}
\item load the midi file rmg1.mid in the track window
\item load the harmonies rmg1.har in the harmony browser. You may minimize the
harmony browser but not close it - it is needed to transpose the generated melody.
\item load the melody generator settings rmg1solo.mel into the Random Melody Generator
\item next, mark the range where the events should be generated. This is
bar 5 to bar 45 in the track window.
\item In the Random Melody Generator window press the generate button,
then the transpose button and finally the play button (the three rightmost
buttons in the toolbar).
\end{itemize}



\subsection{Event List Editor}\label{evtlst}

In this window, you can edit events like text. When the window is opened,
JAZZ++ translates the selected events into a text form. You can edit this
text in the Event List Editor and finally convert it back to midi events.

The Event List Editor allows to filter events. The Editor will show only those
events that are currently enabled in the \helpref{Filter Dialog}{midifilter}
in the Settings menu. As we already know, changing the filter settings
affects most of the other edit functions too, e.g erase will erase only
those events, that are enabled in the Filter Dialog.

The Event List Editor recognizes the following events:

\begin{itemize}
\item {\tt 1:2:030 key 1 4C 90 200}\\
      A note event. 1:2:030 is the timing information, it means bar 1,
      beat 2 plus 30 ticks. The string key indicates that this line
      corresponds to a note on event. The following 1 is the midi channel
      (numbered from 1 to 16), 4C is the octave and pitch, 90
      the velocity and 200 the length of the note.

\item {\tt 1:2:030 kpr 1 4C 90}\\
      A keypressure event. It defines a pressure on channel 1, value 90 for the note 4C.

\item {\tt 1:2:030 ctl 1 10 40}\\
        A controller event. 1 is the midi channel, 10 the controller number (pan) and
        40 the value.

\item {\tt 1:2:030 prg 1 5}\\
        A program change event, here program 5 on channel 1)

\item {\tt 1:2:030 cpr 1 50}\\
        A channel pressure event, here value 50 on channel 1

\item {\tt 1:2:030 pit 1 -4074}\\
        A pitchbend event, valid values are in -8191 .. 8191

\item {\tt 1:2:030 sys f0 11 22 33 44 f7}\\
        A system exclusive message, hex values including the start-
        (f0) and stop- (f7) byte.

\item {\tt 1:2:030 txt here is some text}\\
        A text event. All text up to the end of the line is text.

\item {\tt 1:2:030 cpy copyright text}\\
        A copyright event. All text up to the end of the line is text.

\item {\tt 1:2:030 trn copyright text}\\
        A trackname event. All text up to the end of the line is text.

\item {\tt 1:2:030 spd 120}\\
        A set tempo event. The value is in beats per minute.

\item {\tt 1:2:030 tsi 7/8}\\
        A time signature event. The denomiator must be a power of 2, e.g one
        of 2,4,8,16,32.

\item {\tt 1:2:030 ksi maj 3}\\
        A key signature event. The string 'maj' indicates a major scale,
        'min' would mean a minor scale. The number is the number of
        sharps in the scale.

\end{itemize}



\subsection{Arpeggio Generator}\label{arpeg}

This dialog generates arpeggios from chords that are already recorded.
The generator scans the selected track(s) in the trackwindow. Whenever
3 or more notes sound simultanously, they are replaced by an arpeggio that
consists of the chord notes and their octaves.

You can paint the kind of arpeggio: the x axis represents the time in steps,
the y axis the note of the chord. A value of 5 on the y axis means: 'take
the 5-th note of the chord'. If the chord does not contain 5 notes, the
generator transposes the chord by one octave.

If you check the controller box, you can choose a controller (e.g. panpot).
In this mode, the generator does not scan the selection but generates controller
events.

Take a look at the examples in the song/arpeggio directory.


\subsection{Set Music Copyright}\label{muscopyr}
The copyright notice is put into the midi file as a 'copyright' meta event.
It is intended to be used for copyrighting the music.

\section{Help menu}\label{twhelp}

\begin{itemize}

\item {\em Jazz} loads top-level 'table of contents' of the help-file.

\item {\em Trackwin} loads trackwin chapter.

\item {\em Mouse} shows a dialog that describes mouse actions (what the buttons do
etc.).

\item {\em About} shows a dialog with info about Jazz.

\item {\em Welcome} shows a dialog containing the Jazz 'welcome-message'.

\end{itemize}

\section{Help system}\label{help}
Context specific help is available via the ``help'' button in most dialogs.

There is also a help menu in both main windows.

The help buttons and some of the help-submenus requires access to
the wxWin help system. The help system acts a little different dependant
on the system platform:

\begin{itemize}
\item {\em Unix}. A copy of the {\em wxhelp} utility program as well as
the file {\em jazz.xlp} has to be present in a directory listed by the
shell's binary search path (e.g. /usr/local/bin). If the wxhelp program
cannot be found there will be a timeout message
after a while (repeated two times).

Please observe that the jazz.xlp file does not support graphics.
In order to view the help-file images on Unix (can be useful)
please use the html- or postscript-versions of the documentation. The html
files can be viewed using a standard web browser (like 'netscape'). The
postscript file can be viewed using 'ghostview' or it can be produced on paper
using a postcript-capable printer.

\item {\em Windows}. The native help system is used, which requires that the
help file is present in 'Windows help' format ({\em jazz.hlp}). This file is
installed automatically by the jazz install-utility.
\end{itemize}


%%%%%%%%%%%%%%%%%%%%%%%%%%%%%%%%%%%%%%%%%%%%%%%%%%%%%%%%%%%%%%%%%%%%%%%%%

\chapter{Piano window}\label{pianowin}

The pianowindow shows Note-On events of a single Track.
Displayed track and bar is selected by clicking with the right
mouse button into the trackwindow.

\xlpignore{\begin{figure}%
\latexonly{\image{13cm;0cm}{pianowin.ps}}%
\htmlonly{\image{0cm;0cm}{pianowin.png}}%
\winhelponly{\image{0cm;0cm}{pianowin.bmp}}%
\caption{Piano Window}%
\end{figure}}



\section{Toolbar}\label{pwtoolbar}

\xlpignore{\begin{figure}%
\latexonly{\image{13cm;0cm}{pwtoolb.ps}}%
\htmlonly{\image{0cm;0cm}{pwtoolb.png}}%
\winhelponly{\image{0cm;0cm}{pwtoolb.bmp}}%
\caption{Pianowin toolbar}%
\end{figure}}

Buttons to set the mode of the left mouse-button:
\begin{itemize}
\item {\em select area} (by dragging a rectangle)
\item {\em change length} (by dragging the 'tail' of an event)
\item {\em show dialog} (\helpref{edit existing event or create a new one}{pweventedit})
\item {\em cut/paste} (great for fast note-by-note editing)
\end{itemize}

Note that these actions can also be done by using other mouse buttons or
by combining shift- and ctrl-keys with mouse buttons
(see \helpref{Mouse actions}{mouseact}).


Buttons to set the most common \helpref{snap}{snap} values (pasted events
will be placed at nearest 'snap' value, similar to a 'grid' in a graphics
paint system):
\begin{itemize}
\item set snap to 1/8 note
\item set snap to 1/12 note
\item set snap to 1/16 note
\item set snap to 1/24 note
\end{itemize}

Buttons to manipulate events:
\begin{itemize}
\item \helpref{Cut}{pwcutpaste}
\item \helpref{Delete}{pwcutpaste}
\item \helpref{Quantize}{quantize}
\item \helpref{Shift left}{pwshift}
\item \helpref{Shift right}{pwshift}
\item Enable {\em multi track mode} as described \helpref{here}{pwevents}
\end{itemize}

Buttons for other operations:
\begin{itemize}
\item \helpref{Undo}{undo}
\item \helpref{Redo}{undo}
\item \helpref{Panic (MIDI reset)}{panic}
\item \helpref{Help}{help}
\end{itemize}



\section{Mouse actions}\label{mouseact}

For a fast reference to mouse-actions, please invoke the {\em Help-$>$Mouse}
menu entry.

{\bf Mouse actions in the event area}

\begin{itemize}

\item {\em Select events}. When in 'select area' mode
(see \helpref{toolbar}{pwtoolbar}), events can
be selected by dragging a rectangle with the left mouse button.
A selection is continued by dragging while holding the shiftkey down. (To
continue a selection outside the current view, you need to scroll using the
window scrollbars before doing shiftkey + left-button again.)

\item {\em Adjust note length by dragging}. Click on the event and drag
the 'tail' while holding mouse button down. This can be done with the
left mouse button when in 'change length' mode
(see \helpref{toolbar}{pwtoolbar}), or by using the right mouse button.

\item {\em Show event dialog}. Opens a dialog box to
modify the values of the event
or to create a new event. This can be done with the left mouse button when in
'show dialog' mode
(see \helpref{toolbar}{pwtoolbar}), or by using ctrl+middle (if
available).

\item {\em Cut / paste events}. Clicking on an event cuts the event to the
internal cut/paste buffer, clicking on an empty position pastes the
buffer contents into the pianowin. This can be done with the left mouse
button when in 'cut/paste' mode
(see \helpref{toolbar}{pwtoolbar}), with the middle mouse button (if
available) or with shift+ctrl+left.

\item {\em Copy events}. Copy an event into the internal cut/paste buffer. Do
this by pressing shift+ctrl+right or shift+middle (if available).

\item {\em Change velocity of event}. Holding the controlkey while pressing the
left or right button on an event increments/decrements the velocity. The
velocity value is displayed in the left end of the top line (just
above the piano keyboard).

\item {\em Change track in 'single track mode'}. You may change the track by pressing
shift + cursor keys.

\item {\em Change track in 'multiple track mode'}. In this mode you can change track
by pressing the right mouse button on an event from another track
(painted with grey color).

\item {\em Listen to pitch}. Do this by pressing shift+right.

\end{itemize}



{\bf Mouse actions in the top line}

\begin{itemize}
\item {\em Start/stop play}. Click the left mouse button on the top line to
start/stop playing (record is not available in the pianowin).
\item {\em Start cycle play}. Select some events (some bars) and press
shift+left-button. If there is no selection active, 4 bars are cycled.
\end{itemize}


\section{Edit/create single events}\label{pweventedit}

\xlpignore{\begin{figure}%
\latexonly{\image{13cm;0cm}{eventdlg.ps}}%
\htmlonly{\image{0cm;0cm}{eventdlg.png}}%
\winhelponly{\image{0cm;0cm}{eventdlg.bmp}}%
\caption{Example of event editing dialog}%
\end{figure}}

There are a number of dialogs that allow for detailed editing of single events.
The dialogs are invoked by clicking left mouse button in the pianowin when
in {\em show dialog} mode (see \helpref{toolbar}{pwtoolbar}). Events that
can be created or edited are {\em Note On}, {\em Controller},
{\em Program Change}, {\em Set Tempo} and {\em System Exclusive}.


\section{Edit menu}

Some entries work the same way as discussed in the \helpref{edit menu}{twedit}
of the trackwindow. Only the entries that is special for the pianowindow is
described here.

\subsection{Cut, Copy, Paste}\label{pwcutpaste}
Copies selected events into an internal buffer. Events can be pasted to an
empty space.

Cut is like Copy + delete.

\subsection{Shift}\label{pwshift}

Moves selected events to the left or right. The amount is snaps, which can be
adjusted in the settings-menu.

\subsection{Quantize}
Quantizes selected events to the current snap-setting.

\subsection{Left - Right, Up - Down}

Exchanges the position of selected events. Try it.

\section{Settings menu}

Some entries work the same way as discussed in the
\helpref{settings menu}{twsettings} of the trackwindow.

\subsection{Events}\label{pwevents}

\xlpignore{\begin{figure}%
\latexonly{\image{13cm;0cm}{pwevents.ps}}%
\htmlonly{\image{0cm;0cm}{pwevents.png}}%
\winhelponly{\image{0cm;0cm}{pwevents.bmp}}%
\caption{Pianowin show-events dialog}%
\end{figure}}

You can specify what events are to be shown in the pianowin.

\begin{itemize}
\item {\em NoteOn, Controller, Program, Tempo, SysEx}. These events can be
shown and
edited in the pianowin. If you select exactly one, the piano at the left
hand side will be replaced by a list of controller- or program names.
\item {\em Show drumnames}. Switch it off if you prefer to see the piano
keyboard instead of drumnames.
\item {\em Show events from all tracks}. Displays all events
('multi track mode'). Events from active track are shown
red, the others grey. Clicking with the right mouse button on a grey event
changes the actual track.
\item {\em Show harmonies from harmonybrowser}. Displays the chords and
scales from the harmony browser. To see the chords, you must first open the
harmonybrowser and define a chord sequence. Then you have to mark some bars
in the trackwin to show, where the chords should go. Now the pianowin shows
notes from chords in light blue and notes from the scale in light green.

\end{itemize}

\subsection{Snap}\label{snap}

The start time of pasted events will be quantized to this value. This value is also
used in the quantize- and shift dialogs in pianowin and trackwin.


\section{Misc menu}

\subsection{Edit Pitch and others}

\xlpignore{\begin{figure}%
\latexonly{\image{13cm;0cm}{pitched.ps}}%
\htmlonly{\image{0cm;0cm}{pitched.png}}%
\winhelponly{\image{0cm;0cm}{pitched.bmp}}%
\caption{Graphical editing in Pianowin}%
\end{figure}}


Opens a painting area where you can paint values of the specified type (pitch, controller, ...).
Painting with mouse is done analogous to the \helpref{random rhythm generator}{randrhy}, make sure
you use the right button to set values to zero because you can't paint an exact zero because
of limited screen resolution. Use the Apply button to make your changes take effect.

Also 'tempo' change events can be painted. Please note that tempo changes
in the middle of the song may not work if audio/midi synchronization is
enabled (see \helpref{Global Audio dialog}{globalsettings}).

\subsection{Guitar board}

Opens a window displaying the frets of a guitar. Moving the mouse over the guitar board shows
the names of the notes and moves the black bar in the pianowin. Same works the other way around,
moving the mouse in the pianowin shows the notes on the guitar board.

Also, the actual contents of the pianowin cut / paste buffer is shown (this may be set from the
harmony browser too, so you can see how chords could be played).

Clicking with the left button adds events to the pianowin buffer so you can paste them into
your song. In chord mode (settings dialog), if you add several events to the buffer, they all
start at the same time, otherwise they are generated one after the other (length is the
current snap setting).


\section{Help menu}

\begin{itemize}

\item {\em Pianowin} - loads the pianowin chapter of the help file.

\item {\em Mouse} - displays a dialog that describes mouse actions (what the buttons
do etc.).

\end{itemize}

For more info about help see \helpref{'Help System'}{help} in the {\em Track Window} chapter.

%%%%%%%%%%%%%%%%%%%%%%%%%%%%%%%%%%%%%%%%%%%%%%%%%%%%%%%%%%%%%%%%%%%%%%%%%

\chapter{Audio}\label{audio}

\section{Audio Overview}\label{audiointro}

Audio signal processing is a very CPU and memory intensive task. For
satisfying results we recommend a pentium processor with 32 MB RAM or
up.

The audio part of JAZZ++ is organized similar to a sampler or a
drum machine. That is, you can assign a sound sample to every midi key.
This is different to a usual synthesizer where you assign a sound
to a midi channel and the note only changes the pitch. So an audio track
in JAZZ++ is very much like a drum track. On a drum track, every
note plays another instrument, on an audio track every note plays
another sample.

You assign samples to midi keys using the \helpref{Sample Settings
Dialog}{samplesettings}. JAZZ++ then loads the samples into memory. You
may use samples with different sampling rates and channels (mono,
stereo) in one song, JAZZ++ converts all samples into the format you
specify in the \helpref{Global Settings Dialog}{globalsettings}.
The collection of samples you define in the Sample Settings Dialog
is called a {\em Sample Set} in JAZZ++.

After defining a sample set, you should define an audio track. You do
this in the track window by opening the \helpref{Track Dialog}{trackdlg}
and check the entry called Audio Track. Audio tracks behave the same
way like drum tracks or other midi tracks. E.g you can place individual
notes on the track, copy, delete or replicate bars etc. Every note
on the audio track corresponds to a sample you just defined.

\xlpignore{\begin{figure}%
\latexonly{\image{13cm;0cm}{audiotrk.ps}}%
\htmlonly{\image{0cm;0cm}{audiotrk.png}}%
\winhelponly{\image{0cm;0cm}{audiotrk.bmp}}%
\caption{Audio track in pianowin}%
\end{figure}}

Next open the piano window on that track by clicking with the right
mouse button on the audio track you just created. When you scroll up
and down you should see the labels of the samples on the left hand side
of the window (the labels are assigned in the Sample Settings Dialog).
When you click on one of the labels, you should hear the sound of
that sample.

Next select the 1/16 note in the toolbar to be able to place some notes
on the track. Then set a few notes. When placing a note, you will hear
the sample too. Also you may notice, that the length of the note is
not 1/16 but is as long as the sample is according to the current midi
speed. When you change the midi speed in the trackwin, you will see that
the note length is adjusted to the new speed (by the way: this way you
can easily adjust the midi speed to a drum loop sample).

Now you are ready to listen to your new song. Go back to the track
window and click on the start play toolbar button. And voila, you will
hear your new audio enabled song.


\section{Audio Recording}

Audio record works much the same way like midi record. Select some bars
on an audio track where you want to record an audio signal. Then press
the record button. More on record / play you can read
\helpref{here}{recplay}.  While recording you may not hear other samples
depending on the capabilities of your sound card. Most sound cards do not
support simultaneous record and playback.

When you have finished recording, stop the operation by clicking on the
play or record
button again. JAZZ++ will then ask you, if you want to keep the recording.
When you answer yes, you will get the file selector dialog to specify
the audio file for your recording. The recording will be saved in the
Microsoft .wav format. Also JAZZ++ automatically assigns a midi note to
your new sample and places a note on the recorded audio track. If you
don't like your recording and you do it again, you should save it under
the same file name as the old recording. Then JAZZ++ will reuse the same
midi note and not assign a new one.

Please note: when you do audio recording, JAZZ++ needs much memory for the
playback too. So you should start playback a few bars before the
selection and stop playback a few bars after it. Please do not playback
the whole song, when you only want to record a few bars.


\section{Audio Menu}

\subsection{Global Settings}\label{globalsettings}

\xlpignore{\begin{figure}%
\latexonly{\image{5cm;0cm}{audioset.ps}}%
\htmlonly{\image{0cm;0cm}{audioset.png}}%
\winhelponly{\image{0cm;0cm}{audioset.bmp}}%
\caption{Global audio settings dialog}%
\end{figure}}

Here you specify the internal resolution for the audio processing of
JAZZ++.  This means, that all samples are internally converted to the
parameters you specify here. If you set the sampling rate to 22050 here
and load a sample that was sampled with 11025 samples per second, JAZZ++
will convert it to 22050 samples per second on loading. Specifying
low quality here (low sampling rate and mono) will dramatically reduce
the CPU and memory requirements needed for processing. Of course this
will reduce sound quality too.

\begin{itemize}

\item {\em Enable Audio}. If you leave this unchecked, the whole audio system
will be disabled. This is useful, if you do not use audio in a midi-only
song and want to save CPU and memory.

\item {\em Sample Freq (or Sampling Rate)}. The number of samples per second and
per channel (if stereo). 44100 is CD quality, 8000 is telephone sound.

\item {\em Stereo}. Switch this off if all your samples are mono anyway or to
save memory.

\item {\em Software Midi/Audio Sync}. You may try to switch this off. If midi
and audio keep in time, then your sound card synchronizes midi and audio
itself and you should leave this unchecked, because midi-audio synchronization
is CPU intensive too. Most cards do not synchronize midi and audio, so
in most cases this should be checked. Please note that midi tempo changes
in the middle of a song will not work if Midi/Audio Sync is enabled.

\end{itemize}



\subsection{Sample Settings}\label{samplesettings}

\xlpignore{\begin{figure}%
\latexonly{\image{7cm;0cm}{sampset.ps}}%
\htmlonly{\image{0cm;0cm}{sampset.png}}%
\winhelponly{\image{0cm;0cm}{sampset.bmp}}%
\caption{Sample settings dialog}%
\end{figure}}

In this dialog you define the midi keys for your samples. Please read
the \helpref{Introduction}{audiointro} for more information on midi
keys for samples.

You add a new sample to the list by selecting an list entry first
and then press the add button. In the file selector you choose the
sample file you want to play in your song. The only file format supported
currently is the Microsoft .wav format. After loading the file you should
enter a label (a name for the sample) in the label text entry near the
bottom. This label will be shown in the piano window when editing an
audio track.

The list shows the available midi keys at the left. Low numbers at the
top of the list stand for very low midi keys and high numbers for high
midi keys. To avoid scrolling in the piano window you probably want to
load the samples somewhere in the middle of the list.

To remove a sample from the list select that entry and press the Clear
button. To listen to a sample press the Play button, to interrupt
listening press it again. The Edit button invokes the \helpref{Sample
Editor}{samplewin}.

The Sample Settings dialogs also allows to adjust volume, pan and pitch
of a sample. In opposite to the Sample Editor, these changes are not
permanent, e.g if you reduce the volume you can increase it later again
with no loss of quality. JAZZ++ applies these settings every time a sample
is loaded.

After you have defined a set of samples you want to use, you may
save these settings into a .spl file. Close the dialog and select Save Set from the
Audio menu in the track window. Please notice, that these files contain all
your settings from the Sample Settings Dialog, but they do not contain
the samples themselves. The .spl files only contains the filenames of
the samples, so if you rename or delete samples, JAZZ++ will not be able
to find those samples again.

\section{Sample Editor}\label{samplewin}

\xlpignore{\begin{figure}%
\latexonly{\image{13cm;0cm}{sampedit.ps}}%
\htmlonly{\image{0cm;0cm}{sampedit.png}}%
\winhelponly{\image{0cm;0cm}{sampedit.bmp}}%
\caption{Sample editor}%
\end{figure}}

The sample editor may be invoked from the
\helpref{Sample Settings Dialog}{samplesettings} or from the piano window.
In the piano window, the editor is invoked when you select the dialog
toolbar button on an audio track and click on a note.

The sample editor shows the waveform of the audio signal and allows
you to modify it. There are two scrollbars, the upper one controls the
zoom and the lower one the view position.

For some operations (like cut) you must select a range of samples, you
can do this by dragging the mouse. Similar to the track window or piano
window you may continue a selection by dragging with the shift key down.
Clicking again clears the selection. Other operations (like paste) need
an insertion point. This is set by clicking into the sample area when
there is no selection. The insertion point is displayed as a vertical
red line.

A more convenient way to zoom into a part of a sample is to select
the range of interest and then click the magnify glass button.

Pressing the play button will play the whole sample by default. If there
is a selection, only the selected samples will played. If there is
an insertion point, play will start at that point.



\subsection{Loading and Saving Samples}

After you have modified your sample you may save it. Please note that
saving will make the settings in the \helpref{Sample Settings
Dialog}{samplesettings} permanent. That is, if you have set the volume
e.g. to a value of 50 in the sample settings dialog, the sample will be
saved with this volume and you cannot make it louder again without loss
of quality. To avoid this, you can set the volume to maximum in the
sample settings dialog and leave the pan and pitch sliders neutral before
starting to edit. After saving your sample you can change volume etc again.

You may think of the sample settings dialog as some kind of mixer
whereas the sample editor really changes the sound source itself.

\subsection{Cut/Copy/Paste}

\begin{itemize}

\item {\em Cut}. Cuts out a piece of the sample into an internal buffer. Needs
some samples to be selected.

\item {\em Copy}. Copies some samples into the internal buffer. If no range is
selected, the whole sample is copied.

\item {\em Paste}. Inserts the samples from the internal buffer to the sample.
If there is a selection, the selection is replaced by the buffer. If there
is an insertion point, the buffer is inserted at that position.

\item {\em Paste Merge}. Adds the buffer at the insertion point. The samples
from the buffer are mixed with the original.

\end{itemize}



\subsection{Silence}

You have to select a range before using this function.
\begin{itemize}
\item {\em Replace} - the selection will be replaced with silence.
\item {\em Insert} - will insert silence at the beginning of the selection,
the length of the inserted silence is equal to the length of the selection.
\item {\em Append} - will append silence at the end of the selection. The
length is equal to the length of the selection.
\end{itemize}


\subsection{Invert Phase}
Inverts the phase of the left or right channel


\subsection{Maximize Volume}
This maximizes the volume of the whole sample. It is made as loud as possible
without distortion.


\subsection{Painters}\label{audiopainters}

After selecting a painter command in the menu bar, a paintable area is
displayed at the bottom of the window. In this area you may paint some
kind of curve. Painting works the same way as in the \helpref{Random
Rhythm Generator}{randrhy} and is described there in more detail. After
painting you can press the green accept-button in the toolbar or the red
crossed cancel button. The operation is applied for the visible range
of samples. So if you don't want to modify the whole sample, zoom to
the range of interest before invoking the painter.

\begin{itemize}
\item {\em Volume} - lets you paint some kind of envelope. The painter lets you
paint the relative volume, that means no painting will not change the
volume, painting above the zero line will increase the volume, painting
below the zero line will decrease it.
\item {\em Panpot} - changes volume between left and right channel
\item {\em Pitch} - changes the pitch of the sample. This can operate on the
whole sample only. The range of transposing can be adjusted in the
Settings menu.
\item {\em Filter} - a dynamic filter, e.g you can paint the cut off frequency.
The filter settings can be adjusted in the settings menu.
\end{itemize}

% how to use, Pitch Painter Settings
% volume pan pitch filter


\subsection{Distortion}

This adds a distortion effect to your sample. In the dialog you can paint
a distortion curve. When pressing Ok, every sample is mapped to a new
value according to the curve. JAZZ++ looks up the original sample value at the
horizontal axis and replaces the sample value with the height value of
the curve at that point. Because painting smooth curves is difficult, there
are some presets, try them out!



\subsection{Reverb}

The reverb algorithm of JAZZ++ is produces a high quality reverb effect.
Parameters of the Reverb Dialog:

\begin{itemize}
\item {\em Room Size} - a value corresponding to the time until the first
echo of the reverb can be heard. If you think of a big room producing
the reverb, this corresponds to the size of the room.
\item {\em Brightness} - controls the frequencies that are echoed from
the rooms walls.
\item {\em Length} - controls how many echos are computed. This controls the
depth of the reverb.
\item {\em Effect Volume} - the final output mixer, you set the volume of the
reverb-only signal here which is mixed to the original then.
\end{itemize}



\subsection{Echo}

In this dialog you can add one or more echos to your sample.
\begin{itemize}
\item {\em Number of echos} - controls how many echos are generated
\item {\em Delay} - the time in milli seconds between the echos
\item {\em Amplitude} - the relative Amplitude in \%. A value of e.g. 20 means,
that the first echo has a volume of 20\% of the original, the second
echo has a volume of 20\% of the first echo and so on.
\item {\em Random delay}. If this is checked, the time between the echos is
chosen randomly. The Delay-Value (see above) is taken as a basis
for the random generator.
\end{itemize}




\subsection{Chorus}

This is some kind of universal effect which can produce chorus-like effects.
The operation is quite simple: JAZZ++ takes a copy of the original sample,
then modifies it according to the settings in the dialog and finally
mixes the modified copy back to the original.

\begin{itemize}

\item {\em Delay} - delays the signal for the number of seconds/100 specified
here (e.g. a value of 20 means 200 millisec delay).

\item {\em Pitch Modulation Frequency} - changes the pitch of the signal up
and down controlled by a sine wave LFO (Low Frequency Oscillator), this is
the heart of most chorus effects.
The modulation frequency is set in Hertz/10 (that is, a value of 20 means 2
turnarounds per second). The modulation frequency corresponds to the speed
regulator of a chorus effect.

\item {\em Pitch Intensity} - the depth of pitch change.

\item {\em Pan Freq} - changes the pan of the signal to the left and right
controlled by a sine LFO.

\item {\em Stereo Spread} - enhances the stereo effect of the output. You
may even try to convert a mono signal to a stereo signal by setting Delay
to 20 or so, turn pitch intensity and pan spread off and set stereo spread
to maximum.

\item {\em Effect Volume} - the volume of the modified signal in \% of the
original.

\end{itemize}



\subsection{Pitch shifter}\label{pitchshifter}

This operation lets you change the pitch of a sample. Pitch changing
normally changes the length of the sample too (you can think of a tape
machine, where the speed is altered). JAZZ++ however provides algorithms
to change the speed without changing the length of the sample.

With the sliders you specify the amount of transposition in semitones.
If the box {\em Keep Length} is checked, the transposed sample will have the
same length as the original. In this case, JAZZ++ divides the sample
into small pieces and plays these pieces with the new speed as often as
necessary to match the original length. This may lead to an undesired
tremolo effect if the size of the pieces is too big or too small. The
slider named window size lets you adjust the size of the pieces.


\subsection{Time Stretcher}

This is similar to the \helpref{Pitch shifter}{pitchshifter}. You can
change the length of the sample here. Changing the length normally
changes the pitch of the sample too, but if you check the box labeled
{\em Keep Pitch} JAZZ++ will change the length without changing the pitch.

You may specify the new length in seconds (upper sliders) or in terms
of midi speed (lower sliders). With the upper sliders you specify the
delta length, that is how much time you like to add or subtract from
the original length. Eg to make the sample one second shorter, you
would set the first slider to -1.

The midi speed sliders are best explained using an example: Lets assume
you have a drum loop recorded at a speed of 120 beats per minute. To
change the speed (and thus the length) of the sample, to 130 beats per
minute, you set the slider {\em old speed} to 120 and the slider {\em new speed}
to 130.

The meaning of the window slider is described in the \helpref{Pitch
Shifter Dialog}{pitchshifter}.



\subsection{Reverse}

This reverses the output (think of a tape playing backwards). If there
is no selection, the whole sample is reversed.



\subsection{Additive Synthesis}

In this window you can generate new sounds using a noise generator
together with additive sound synthesis generators. You may up to 6
generators to build your sound. For each generator you can paint curves
for envelope, panpot and pitch. The output of all generators is mixed
together to make the final sound.

The toolbar buttons have the following meaning (from left to right):
\begin{itemize}
\item load synth settings. In the song subdirectory you find an
example for a simple drum which maybe a starting point for you.
\item save a synth settings.
\item generate a wave
\item listen to the generated wave (you have to generate first). If you
dont hear anything make sure that you have painted some envelope. Also
you have to enable audio in the audio menu of the track window of course.
\end{itemize}

Near the top of the window you find some checkboxes. The first
4 (Harmonics, Envelope, Panpot and Pitch) switch on and off the
corresponding paint areas.  This is useful for painting on a small screen.
The last one labeld {\it Noise} turns the first generator into a noise
generator when checked.

{\bf Noise generator}

In the left paint area labeled {\it noise filter}  you can paint the
cut off frequency of the noise generator. The slider labeled {\it Midi Key}
is ignored in the noise mode.


{\bf Sound generator}

In sound generator mode the left paint area is labeled {\it harmonics}. You can
paint the volumes for each harmonic. Note that the 10 leftmost harmonics
determine much of the final sound (you have to make the window big to
paint these). Harmonics with even numbers sound warmer than harmonics
with odd numbers. Most sounds will only have a few harmonics with
high volume and especially harmonics with high numbers will have low
volume. In this mode the midi key is the base frequency from which the
harmonics are calculated.

{\bf working with multiple generators}

With the slider labeled {\it Number of Synths} you can use up to 6 generators
simultanously. Only the first one may be a noise generator, all others will
be in additive synthesis mode. There is much room for experimenting with
multiple generators, e.g you can do some kind of morphing / crossfading with
the envelope painters, or get some chorus like effects when detuning the
oscillators a little bit in the pitch painter.

After you have made some sound make sure to add some reverb or chorus.



\subsection{Filter}

This adds high quality filters to your effects gallery. You may choose
one of the following filter types:

\begin{itemize}
\item {\em Low Pass} - lets pass low frequencies, the volume of high
frequencies is reduced (some kind of bass booster).
Frequencies below the corner frequency will become louder, frequencies above
will become more quiet. The bandwidth value is ignored.

\item {\em High Pass} - the other way round, the volume of low frequencies is
reduced and high frequencies can pass the filter (some kind of treble booster).
Frequencies above the corner frequency will become louder, frequencies below
will become more quiet. The bandwidth value is ignored.

\item {\em Band Pass} - increases the volume of middle frequencies. The corner
frequency is the frequency with maximum volume. The bandwidth is a value
in \% of the corner frequency and specifies, what frequencies above
and below the corner frequency may pass the filter.

\item {\em Band Stop} - decreases the volume of middle frequencies.
\end{itemize}

The feedback value controls the gradient of the filter. In case of a low
pass filter for example, it controls how much the volume above the
corner frequency is reduced. A value of 1 or 2 is sufficient in most
cases.

\subsection{Settings}

\begin{itemize}

\item {\em Pitch Painter} opens a dialog to control the range of the
pitch paint area (see \helpref{Painters}{audiopainters}).

\item {\em Filter Painter} opens a dialog to control the type of filter that
is applied to the filter painter area (see \helpref{Painters}{audiopainters}).

\item {\em View Settings} opens a dialog to control view parameters of the
sample editor window. The 'show midi time' checkbox toggles between
showing time as counts or real time.

\end{itemize}

%%%%%%%%%%%%%%%%%%%%%%%%%%%%%%%%%%%%%%%%%%%%%%%%%%%%%%%%%%%%%%%%%%%%%%%%%

\chapter{Examples}\label{examples}

\section{How to build a new midi song}\label{exmidisong}

Supposing you start with a blank song in the trackwin, these are the steps
you need to do in order to record some midi tracks and make your own midi song.

\begin{enumerate}

\item {\em Turn on the metronome}. Push the \helpref{metronome button}{twtoolbar}.
This will help you keep in time when playing.
Later, when you have recorded some tracks, you may want to turn the metronome
off again because the other recorded tracks will support you with the necessary
beat information.

\item {\em Select counts per bar etc.}. By default, a blank song will have
4 counts per bar with each count corresponding to a quarter note. To change
this you need to invoke the \helpref{Misc-$>$Meterchange}{meterchange}
menu entry.

\item {\em Select song speed}. You probably have a feeling about how fast
your song should be. You can try this right away by just pressing the
\helpref{play button}{twtoolbar}. You will then hear the beats of the
metronome. Change the speed to your preference by clicking into the
\helpref{speed field}{speed}.

\item {\em Define track defaults}\label{rectrackdef}. Before you can enter music
(events) into
a track, you need to define {\em track name}, {\em instrument to be played}
and {\em midi channel number} for that track. Do this by clicking on the
\helpref{track name field}{trackdlg}, starting on the {\em second} track (will
invoke the Track Settings dialog). {\em Never use the first track for
recording}, as this is defined to be the {\em master track} (holding events
controlling the whole song like speed, meter change etc.).

\begin{itemize}

\item {\em Track name}. Enter a track-name that describes what you play on the
track, not just the instrument
name (you may want to play some other voice using the same instrument into
another track at a later stage).

\item {\em Instrument}. Select an instrument from the instrument list. The instruments with
one number in front of the name (e.g. '5 Rhodes Piano') are the 'General MIDI' (GM)
instruments, the others are extensions to the GM standard and may not be
supported by your synth/sound-card.

If this track is to play the same
channel and the same instrument as a previously defined track, set the
instrument to 'None' (default).

If your track is a drum track, the instrument names are a bit misleading.
To get started we recommend that you select instrument number 1
(labelled '1 Piano 1'). This will give you the 'standard drum kit' on a
GM compatible synth/sound-card. You must study the instrument list in your
synth documentation to see what other drumkits are available.

\item {\em MIDI channel}. If the track is to play a new instrument, choose a
midi channel that
has not been used by another track. Start at channel 1 and increment by one
for each new instrument (track). If your track is to play {\em drums}, select
channel 10 (predefined by GM to be the drum channel).

\end{itemize}

\item {\em Set output midi channel on your keyboard}\label{recsetkbchan}.
If your keyboard is
capable of sending on different midi channels, you should now set the
keyboard output channel according to the midi channel of the track you are
about to record. If your keyboard can only send on one midi channel
(e.g. channel 1), you must use a slightly different algorithm
(described \helpref{below}{onechannel}).

Note that you can also 'record' without a keyboard by manually entering
events into the \helpref{pianowin}{pianowin} using the mouse.

\item {\em Select some bars}. \helpref{Select}{select} the area of your track
where you want to record. You can also select the whole track by
clicking in the leftmost field of the track in trackwin (the
\helpref{track/channel display column}{trackdefs}).

\item {\em Start recording}. Start recording by clicking at the
\helpref{record button}{twtoolbar}. JAZZ++ will then start to play one bar
ahead of the selected area to give you a count-in. You can also start
recording by clicking directly at a bar in the top-line of the event area
(see section \helpref{Record and play}{recplay}). While recording you will
hear the metronome (if enabled), any previously recorded tracks plus your own
keyboard playing.

\item {\em Stop recording}. Stop recording by clicking the
\helpref{record button}{twtoolbar} again (or use one of the other 'stop-play'
methods). If you played any notes, you will be asked if you want
to keep the recorded events. If you decide to keep them, press 'Ok'. You can
then listen to the recording. If you are not satisfied you can press the
\helpref{Undo button}{twtoolbar} to take it away and then do a new recording
(go back to previous item).
You can also do a new recording by pressing right-button on the top-line.
The old recording will then be silent, and any new recording will replace the
old one.

\item {\em Post process the recording}. You may want to fix something in your
recording afterwards. You may want to \helpref{quantize}{quantize} or
\helpref{transpose}{transpose} the recording (menu entry {\em Edit-$>$Transpose}),
or you may want to fix notes that are wrong using the mouse in the
\helpref{pianowin}{pianowin}.

\end{enumerate}

You can record a new track by going back to the item
{\em 'Define track defaults'}. You can also record more
music into an existing track by going back to item
{\em 'Set output midi channel on your keyboard'}.

At some point you will probably also want to edit other track defaults like
volume, pan, reverb and chorus. You can do this by invoking the
\helpref{Mixer}{mixer} dialog (press the \helpref{mixer button}{twtoolbar}).

If your synth is {\em GS} or {\em XG} compatible, you may also want to do other kinds of
part editing like adjusting {\em sound-color}, {\em vibrato} or {\em envelope},
or change the behavior of controllers like {\em pitch bender} and
{\em modulation wheel}. These things and a lot more are accessible through the
\helpref{Parts menu}{partsmenu}.

\subsection{Recording with a one-channel keyboard}\label{onechannel}

If your keyboard can only send on one midi channel (e.g. channel 1) you must
change the above procedure according to the following:

\begin{itemize}

\item Always record in a track having the default channel set to 1.

\item When finished recording on that track, change the default track number
to something else (a free channel number) by clicking the track-name-field
and adjusting the slider in the "Track Settings" dialog. When you click
'Ok', the default channel and all events on the track will be converted to
the new channel number (this is true if the entry
{\em 'Force channel number onto all events on track'} is checked).

\item If you want to do more recording on a previously recorded track, change
that track back to channel 1 first using the same method.

\end{itemize}

Pressing the \helpref{Panic}{twtoolbar} button between operations will
prevent the synth/sound-card from being confused (sends a 'reset' command
to the synth/sound-card).


\section{Random rhythm and harmonizer example}

This section documents the demo example.

\begin{enumerate}

\item I prepared the tracks for the song and assigned names and channels for
drums, bass, piano and melody. Then I generated the drums with the
random rhythm generator.

\item open the random rhythm generator and load the file demo.rhy. Lets take a
look at the adjustments for the instruments:

The first kick drum defines high probabilities at beat 1/4 and 3/4 -
this makes the basic groove. The probabilities at 3/16 and 9/16 are a
little lower, so the kick drum will not play these beats always. The
velocity ranges from about 70 to 100, so the loudness of the beats varies.

The snare drum plays at 2/4 and 4/4 but probabilities are less than 100\%, so
it will play most of the time but not always. The velocity is about 80 .. 100, so
the snare is loud always, there is not much dynamic.

The third entry, closed hi-hat, plays at 2/8, 4/8, 6/8 and 8/8 - the probabilities
are 100\% and velocity is high - these beats are played always and all of them are
loud.

The ride cymbal plays most of 1/4, 2/4, 3/4 and 4/4 and sometimes on 16/16. Also
it plays at 9/16 - just to make things more interesting.

The next entry, snare drum again, plays with low probability on all 1/16 timings -
so not all of them are played - the probabilities would lead to about every 2nd
1/16 would be played. For every beat that is played, the generator selects a
length from the length field. Valid length values are 1/16, 2/16 and 3/16 where
1/16 is chosen most often because it has the highest probability. When a length
greater 1 is selected, then the generator will not produce a beat on the following
1/16 step, because the previous note has not finished. Anyway, the velocity is
very low compared to the other snare drum entry. Together both snare drum entries
produce a loud and steady beat from the first entry and a lot of low-velocity
fillings from the second entry.

The last entry, closed hi-hat again, is similar. It does not play on the timing positions
where the first closed hi-hat plays its loud and steady beats. It plays at the other
timing position in between with a very wide range of velocity - so some of these
in between beats will be dominant while others will sound like background.


\item To produce the rhythm I selected bar 5 to 20 on the drum track and then
clicked on the 'gen' button. This gives demo1.mid

\item next I recorded 2 bars of a bass line and some chords. I quantized
this and set the velocity to 90. Last I choose repeat copy to duplicate
the material to the full length of 16 bars. This is demo2.mid.

\item next I recorded a melody and did some corrections in the pianowin editor.
This is demo3.mid.

\item last I wanted some more interesting harmonies. To reproduce this step
you may open the harmony browser and load the file demo.har.

The first 4 chords
are from the original recording Am7 and D7. These are step II and V from G major.
So I chose Gj7 as the next 2 chords to establish that scale.

The next 6 chords are a minor II-V-I move, which leads to Bbm7 in the
scale of Ab major - notice, that the V was taken from the harmonic scale.
The last chords are a V-I move in Ab major.


\item to perform the transposition, transpose every track on its own. Mark the
bass track in the trackwin and select transpose in the harmony browser (the button
where the arrow points from the rectangle to the note). Then do the same for
piano and melody. This should result in demo4.mid.

demo4.mid is not perfectly transposed, but it could serve as a starting point and
you could correct some notes in the pianowin editor. The pianowin editor can
show the chords and scales in colors, look at \helpref{Pianowin Menu Events}{pwevents}
for details.

\end{enumerate}


\section{Audio example}\label{exaudio}

This is how you run the supplied audio demo:

\begin{enumerate}

\item Load the midi file named 'audiodem.mid'. You will see there are 3 tracks
defined; drums, bass and an audio track named 'guitar'.

\item Invoke menu entry {\em Audio$->$Load Set} and load file 'audiodem.spl'.
This is the 'sample set definition file' that binds audio sample files to
the different key pitches (see \helpref{Audio Overview}{audiointro}).

\item Check that audio is enabled (invoke menu entry
{\em Audio$->$Global Settings}, see \helpref{Global Settings}{globalsettings}).

\item Now you can play the song (with audio) (see
\helpref{Record and Play}{recplay}).

\end{enumerate}

You may also want to look at how the samples are bound to the key pitches by
invoking
\helpref{Sample Settings dialog}{samplesettings} and scroll the listbox down to
position 50. There are two samples defined (funkguit1.wav and funkguit2.wav).
Play the samples by selecting in the listbox and press the {\em Play} button.
Look at how you can edit the samples in the
\helpref{Sample Editor}{samplewin} by pressing the {\em Edit} button.

At last you should look at the 'guitar' audio track in the
\helpref{Piano Window}{pianowin}. Scroll
up and down until you find the defined samples.


\end{document}

